\subsection{VIII)}

{\large
    \begin{align*}
    F(A, B, C, D) &= \prod(0, 4, 6, 7, 8, 9, 13, 15) \quad x = 1;
    \end{align*}
}

\subsubsection{Obtener la mínima expresión booleana.}

\begin{center}
\begin{tabular}{c|cccc|c}
\rowcolor[HTML]{FFCCC9} 
\cellcolor[HTML]{9AFF99}$M$ & $A$ & $B$ & $C$ & $D$ & \cellcolor[HTML]{FFFFC7}$F(A,B,C,D)$ \\ \hline
0                           & 0   & 0   & 0   & 0   & 0                                    \\
1                           & 0   & 0   & 0   & 1   & $x$                                     \\
2                           & 0   & 0   & 1   & 0   &                                      \\
3                           & 0   & 0   & 1   & 1   &                                      \\
4                           & 0   & 1   & 0   & 0   & 0                                    \\
5                           & 0   & 1   & 0   & 1   &                                      \\
6                           & 0   & 1   & 1   & 0   & 0                                    \\
7                           & 0   & 1   & 1   & 1   & 0                                    \\
8                           & 1   & 0   & 0   & 0   & 0                                    \\
9                           & 1   & 0   & 0   & 1   & 0                                    \\
10                          & 1   & 0   & 1   & 0   &                                      \\
11                          & 1   & 0   & 1   & 1   &                                      \\
12                          & 1   & 1   & 0   & 0   &                                      \\
13                          & 1   & 1   & 0   & 1   & 0                                    \\
14                          & 1   & 1   & 1   & 0   &                                      \\
15                          & 1   & 1   & 1   & 1   & 0                                   
\end{tabular}
\end{center}

\begin{center}
    \begin{karnaugh-map}[4][4][1][$C\qquad D$][$A\qquad B$]
        \maxterms{0,4,6,7,8,9,13,15}
        \terms{1}{$x$}
        \implicant{0}{4}\implicant{7}{6}\implicant{13}{15}\implicant{8}{9}
    \end{karnaugh-map}
\end{center}

\begin{itemize}[nosep]
    \item Primer bloque ($M_0$ y $M_4$) cambia $B$ y quedan sin negar el resto de entradas.
    \item Segundo bloque ($M_7$ y $M_6$) cambia $D$ y hay que negar $B$ y $C$.
    \item Tercer bloque ($M_{13}$ y $M_{15}$) cambia $C$ y hay que negar todas las entradas individualmente.
    \item Cuarto bloque ($M_{8}$ y $M_{9}$) cambia $D$ y hay que negar $A$.
    \item La función simplificada queda:
\end{itemize}

\begin{center}
    \begin{align*}
        F(A,B,C,D) &= (A+C+D)\cdot{}(A+\overline{B}+\overline{C})\cdot{}(\overline{A}+\overline{B}+\overline{D})\cdot{}(\overline{A}+B+C)
    \end{align*}
    \end{center}

    \subsubsection{Simulación}
    \imagen[]{10cm}{./imagenes/ejercicio8Falstad.png}
    \subsubsection{Descripción en Verilog}\begin{lstlisting}
module descripcionVerilogEjercicio8(
    input inputA,
    input inputB,
    input inputC,
    input inputD,
    output outputFunction
    );
	 assign outputFunction = (inputA | inputC | inputD) & (inputA | ~inputB | ~inputC) & (~inputA | ~inputB | ~inputD) & (~inputA | inputB | inputC);


endmodule
 \end{lstlisting}\subsubsection{RTL output}\imagen[Bloque RTL]{15cm}{./imagenes/rtlEjercicio8Verilog.png}\imagen[Implementación con compuertas]{15cm}{./imagenes/rtlCompuertasEjercicio8Verilog.png}\subsubsection{Simulación temporal}\noindent Código del módulo Verilog Test Fixture:\begin{lstlisting}
module verilogTestFixture;

	// Inputs
	reg inputA;
	reg inputB;
	reg inputC;
	reg inputD;

	// Outputs
	wire outputFunction;

	// Instantiate the Unit Under Test (UUT)
	descripcionVerilogEjercicio8 uut (
		.inputA(inputA), 
		.inputB(inputB), 
		.inputC(inputC), 
		.inputD(inputD), 
		.outputFunction(outputFunction)
	);

	initial begin
		// Initialize Inputs
		inputA = 0;
		inputB = 0;
		inputC = 0;
		inputD = 0;

		// Wait 100 ns for global reset to finish
		#100;
      inputA = 0; inputB = 0; inputC = 0; inputD = 1;
		#100;
      inputA = 0; inputB = 0; inputC = 1; inputD = 0;
		#100;
      inputA = 0; inputB = 0; inputC = 1; inputD = 1;
		#100;
      inputA = 0; inputB = 1; inputC = 0; inputD = 0;
		#100;
      inputA = 0; inputB = 1; inputC = 0; inputD = 1;
		#100;
      inputA = 0; inputB = 1; inputC = 1; inputD = 0;
		#100;
      inputA = 0; inputB = 1; inputC = 1; inputD = 1;
		#100;
      inputA = 1; inputB = 0; inputC = 0; inputD = 0;
		#100;
      inputA = 1; inputB = 0; inputC = 0; inputD = 1;
		#100;
      inputA = 1; inputB = 0; inputC = 1; inputD = 0;
		#100;
      inputA = 1; inputB = 0; inputC = 1; inputD = 1;
		#100;
      inputA = 1; inputB = 1; inputC = 0; inputD = 0;
		#100;
      inputA = 1; inputB = 1; inputC = 0; inputD = 1;
		#100;
      inputA = 1; inputB = 1; inputC = 1; inputD = 0;
		#100;
		inputA = 1; inputB = 1; inputC = 1; inputD = 1;

		// Add stimulus here

	end
      
endmodule
 \end{lstlisting}\imagen[Simulación temporal]{15cm}{./imagenes/simulacionTemporalEjercicio8.png}
